\documentclass[11pt]{article}

\usepackage{series}
\usepackage{booktabs}

\hypersetup{
  pdftitle={Finite Microcanonical Shells and Their Sharp Constraint Limit},
  pdfauthor={Peter Nero},
  pdfsubject={Microcanonical measures, coarea formula, finite energy shells, Hamiltonian invariance, and Modal Triplet Theory},
  pdfkeywords={microcanonical ensemble, coarea formula, energy shell, approximate identity, Hamiltonian flow, Modal Triplet Theory},
  hidelinks
}

\newcommand{\Jcal}{\mathcal J}
\newcommand{\Ecal}{\mathcal E}
\newcommand{\Lip}{\operatorname{Lip}}

\title{Finite Microcanonical Shells and Their Sharp Constraint Limit\\
\large Coarea Geometry, Quantitative Bias, and the MTT Source Boundary}

\author{Peter Nero}

\date{July 2026, Version 1}

\begin{document}
\maketitle

\begin{abstract}
The formal microcanonical density
\(\delta(H-E)\) is a surface distribution, not an ordinary phase-space
probability density. This paper gives a finite-shell formulation and states
precisely when its normalized measures approach the sharp microcanonical
ensemble. For a regular energy \(E\), the coarea formula reduces phase-space
integration to the density-of-states functions
\[
A_f(e)=\int_{H^{-1}(e)}
\frac{f}{|\nabla H|}\,\dd\Sigma_e.
\]
If \(\eta_\epsilon\) is a normalized approximate identity in energy, then the
finite-shell expectation is the ratio of two ordinary convolutions,
\[
\frac{(\eta_\epsilon*A_f)(E)}
{(\eta_\epsilon*A_1)(E)},
\]
and converges to \(A_f(E)/A_1(E)\). When the two density-of-states functions
are locally Lipschitz, an explicit first-order error bound separates shell
width, kernel shape, observable variation, and normalization bias. On a
symplectic phase space, every shell whose density depends only on \(H\) is
exactly invariant under the Hamiltonian flow; ergodicity is not needed for
invariance and is not proved here. Critical energies, noncompact shells,
ensemble equivalence, and a physical choice of shell profile require separate
hypotheses. In Modal Triplet Theory (MTT), the result supplies a rigorous
downstream target: selected geometry may emit a finite energy-tolerance
profile, but neither the coarea theorem nor the sharp limit selects that
profile or its width.
\end{abstract}

\section*{Version 1 Revision Note}
\addcontentsline{toc}{section}{Version 1 Revision Note}
\begin{description}
\item[Supersedes] The unversioned April 2026 manuscript
\emph{Classical Constraint Deltas and Microcanonical Shells:
Admissibility-Shell Limits in Modal Triplet Theory}.
\item[Reason] The earlier manuscript duplicated the general coarea theorem
owned by the delta-kernel paper, described every delta as a finite admissible
process idealized to zero width, and assigned microcanonical closure to the
MTT circle without a source theorem. It did not quantify finite-shell bias,
separate invariance from ergodicity, or explain the failures at critical
energies and on noncompact shells.
\item[Resolution] This version imports the standard coarea reduction and
concentrates on normalized microcanonical measures. It proves convergence,
adds an explicit Lipschitz error certificate, proves exact Hamiltonian-flow
invariance for energy-dependent shell densities, and records the critical,
normalization, ensemble, and MTT-source boundaries.
\item[Retained result] A normalized finite energy shell converges to the sharp
microcanonical surface measure at regular energies under finite
density-of-states hypotheses.
\item[Remaining boundary] No theorem here selects the shell profile or width
from MTT geometry, establishes ergodicity, handles singular critical levels in
general, or proves equivalence with a canonical ensemble.
\end{description}

\tableofcontents

\section{What the microcanonical delta means}

Let \((M,\omega)\) be a \(2d\)-dimensional symplectic phase space with
Liouville measure
\[
\dd\mu=\frac{\omega^d}{d!},
\]
and let \(H:M\to\mathbb R\) be a smooth Hamiltonian. The familiar expression
\[
\rho_E(z)
=
\frac{\delta(H(z)-E)}{\Omega(E)}
\]
is shorthand for a probability measure on the level set \(H^{-1}(E)\), when
that surface measure is finite and nonzero.

If \(E\) is a regular value, then \(\nabla H\neq0\) on \(H^{-1}(E)\), and the
coarea formula gives
\[
\Omega(E)
=
\int_{H^{-1}(E)}
\frac{1}{|\nabla H(z)|}\,\dd\Sigma_E(z).
\]
For an observable \(f\), the sharp microcanonical expectation is
\[
\langle f\rangle_E
=
\frac{1}{\Omega(E)}
\int_{H^{-1}(E)}
\frac{f(z)}{|\nabla H(z)|}\,\dd\Sigma_E(z).
\]
The factor \(1/|\nabla H|\) is not optional. It is the normal Jacobian that
converts ambient phase volume into level-set measure.

The delta notation is exact distribution theory. A finite shell is a different
probability model that can approach it. Such a shell can represent finite
energy preparation, detector resolution, thermodynamic coarse graining, or a
mathematical regularization. Those interpretations are not interchangeable
until a physical preparation protocol is specified.

\section{Finite energy-shell measures}

Let \(\eta:\mathbb R\to[0,\infty)\) satisfy
\[
\int_{\mathbb R}\eta(u)\,\dd u=1,
\]
and define
\[
\eta_\epsilon(u)
=
\frac{1}{\epsilon}\eta\!\left(\frac{u}{\epsilon}\right),
\qquad \epsilon>0.
\]
Compactly supported profiles describe a hard finite band with shaped edges.
The normalized Gaussian
\[
\eta_\epsilon(u)
=
\frac{1}{\sqrt{2\pi}\epsilon}
\exp\!\left(-\frac{u^2}{2\epsilon^2}\right)
\]
describes soft energy mismatch. Both are ordinary nonnegative densities in the
energy variable.

Define the unnormalized shell functional
\[
N_\epsilon(f;E)
=
\int_M f(z)\eta_\epsilon(H(z)-E)\,\dd\mu(z)
\]
and its normalizing factor
\[
\Omega_\epsilon(E)
=
N_\epsilon(1;E).
\]
Whenever \(0<\Omega_\epsilon(E)<\infty\), the finite-shell probability measure
is
\[
\dd\nu_{E,\epsilon}(z)
=
\frac{\eta_\epsilon(H(z)-E)}
{\Omega_\epsilon(E)}\,\dd\mu(z).
\]

The normalization must be checked. On a noncompact phase space, even a narrow
energy profile can have infinite phase volume. Confinement, compact energy
bands, or a declared observable support is therefore part of the theorem, not
a cosmetic assumption.

\section{Coarea reduction to one energy variable}

For energies \(e\) in a regular interval around \(E\), define
\[
A_f(e)
=
\int_{H^{-1}(e)}
\frac{f(z)}{|\nabla H(z)|}\,\dd\Sigma_e(z).
\]
The coarea formula gives
\[
N_\epsilon(f;E)
=
\int_{\mathbb R}
\eta_\epsilon(e-E)A_f(e)\,\dd e.
\]
If \(\widetilde\eta(u)=\eta(-u)\), this can be written
\[
N_\epsilon(f;E)
=
(\widetilde\eta_\epsilon*A_f)(E).
\]
For symmetric shell profiles, \(\widetilde\eta=\eta\).

This reduction is the key simplification. Finite-shell convergence is an
ordinary problem for an approximate identity acting on the density-of-states
function \(A_f\). The geometric work is contained in the coarea factor.

\section{Normalized microcanonical convergence}

\begin{theorem}[Finite-shell microcanonical limit]
\label{thm:microcanonical-limit}
Let \(E\) lie in an open interval \(I\) of regular values of \(H\). Let
\(\eta\geq0\) be integrable with unit integral, and assume its rescaled mass
outside every fixed neighborhood of zero tends to zero. Suppose
\(A_1(e)\) and \(A_f(e)\) are finite and continuous at \(E\), with
\[
\Omega(E)=A_1(E)>0.
\]
Assume also that the shell integrals are finite and that contributions from
outside a compact subinterval of \(I\) vanish as \(\epsilon\downarrow0\).
Then
\[
\int_M f\,\dd\nu_{E,\epsilon}
\longrightarrow
\frac{A_f(E)}{A_1(E)}
=
\langle f\rangle_E.
\]
\end{theorem}

\begin{proof}
The coarea reduction gives
\[
N_\epsilon(f;E)\longrightarrow A_f(E),
\qquad
\Omega_\epsilon(E)\longrightarrow A_1(E)
\]
by the approximate-identity property. Since \(A_1(E)>0\), the denominator is
positive and bounded away from zero for sufficiently small \(\epsilon\).
Taking the quotient proves the claim.
\end{proof}

The theorem does not say that the finite shell is physically more fundamental
than the sharp surface. It says that if a physical or mathematical model uses
such finite profiles, their normalized expectations have a controlled sharp
limit.

\section{A quantitative finite-width certificate}

Convergence alone does not tell us whether a chosen finite shell is already a
good approximation. Let
\[
m_1(\eta)
=
\int_{\mathbb R}|u|\eta(u)\,\dd u
\]
be the first absolute moment of the profile.

\begin{theorem}[Lipschitz shell-bias bound]
\label{thm:bias}
Under the hypotheses of Theorem~\ref{thm:microcanonical-limit}, suppose
\(A_f\) and \(A_1\) are Lipschitz on an interval containing all sufficiently
small shells, with constants \(L_f\) and \(L_1\), and suppose
\(m_1(\eta)<\infty\). Write
\[
a=A_f(E),
\qquad
\Omega=A_1(E)>0.
\]
If
\[
\epsilon m_1(\eta)L_1\leq\frac{\Omega}{2},
\]
then
\[
\left|
\int_M f\,\dd\nu_{E,\epsilon}
-\frac{a}{\Omega}
\right|
\leq
\frac{2\epsilon m_1(\eta)L_f}{\Omega}
+
\frac{2|a|\epsilon m_1(\eta)L_1}{\Omega^2}.
\]
\end{theorem}

\begin{proof}
Changing variables \(e=E+\epsilon u\) gives
\[
|N_\epsilon(f;E)-a|
\leq
\epsilon m_1(\eta)L_f
\]
and
\[
|\Omega_\epsilon(E)-\Omega|
\leq
\epsilon m_1(\eta)L_1.
\]
The assumed inequality implies
\(\Omega_\epsilon(E)\geq\Omega/2\). Therefore
\begin{align*}
\left|
\frac{N_\epsilon(f;E)}{\Omega_\epsilon(E)}
-\frac{a}{\Omega}
\right|
&\leq
\frac{|N_\epsilon(f;E)-a|}{\Omega_\epsilon(E)}
+
\frac{|a|\,|\Omega_\epsilon(E)-\Omega|}
{\Omega_\epsilon(E)\Omega}\\
&\leq
\frac{2\epsilon m_1(\eta)L_f}{\Omega}
+
\frac{2|a|\epsilon m_1(\eta)L_1}{\Omega^2}.
\end{align*}
\end{proof}

The two terms have different meanings. The first measures variation of the
observable-weighted density of states across the shell. The second measures
normalization drift. A small width alone is not enough if the density of states
changes rapidly near \(E\).

For an even profile with a finite second moment and twice differentiable
density-of-states functions, the first signed moment vanishes and a
second-order bound can be obtained. That improvement depends on stronger
regularity and is not used here.

\section{Exact invariance under Hamiltonian flow}

Let \(\Phi_t\) be the Hamiltonian flow of \(H\). Liouville's theorem gives
\[
(\Phi_t)_\#\mu=\mu,
\]
and conservation of energy gives
\[
H\circ\Phi_t=H.
\]

\begin{theorem}[Hamiltonian invariance of every energy-profile shell]
\label{thm:invariance}
Let \(w:\mathbb R\to[0,\infty)\) be measurable and assume
\[
0<\int_M w(H(z))\,\dd\mu(z)<\infty.
\]
Then the probability measure
\[
\dd\nu_w(z)
=
\frac{w(H(z))}
{\int_Mw(H)\,\dd\mu}\,\dd\mu(z)
\]
is invariant under \(\Phi_t\). In particular, every finite shell
\(\nu_{E,\epsilon}\) is invariant whenever it is normalizable.
\end{theorem}

\begin{proof}
For every bounded measurable \(f\),
\begin{align*}
\int_M f\circ\Phi_t\,w(H)\,\dd\mu
&=
\int_M f\circ\Phi_t\,w(H\circ\Phi_t)\,\dd\mu\\
&=
\int_M (fw(H))\circ\Phi_t\,\dd\mu\\
&=
\int_M fw(H)\,\dd\mu.
\end{align*}
Normalization gives invariance of \(\nu_w\).
\end{proof}

Invariance is not ergodicity. A flow can preserve a microcanonical measure
while decomposing its energy surface into many invariant components.
Equilibrium interpretations based on time averages require ergodicity,
mixing, typicality, or another statistical argument not supplied by the delta
or shell construction.

\section{Worked kinetic-energy example}

For one momentum coordinate \(p\in\mathbb R\),
\[
H(p)=\frac{p^2}{2m}.
\]
At \(E>0\), the energy level consists of
\[
p_\pm=\pm\sqrt{2mE}.
\]
Since \(H'(p)=p/m\), the sharp constraint acts on a test function as
\[
\int_{\mathbb R}
f(p)\delta\!\left(\frac{p^2}{2m}-E\right)\,\dd p
=
\frac{m}{\sqrt{2mE}}
\left[
f(\sqrt{2mE})+f(-\sqrt{2mE})
\right].
\]
The density of states is
\[
\Omega(E)
=
\sqrt{\frac{2m}{E}},
\qquad E>0.
\]
The normalized sharp expectation is therefore
\[
\langle f\rangle_E
=
\frac12
\left[
f(\sqrt{2mE})+f(-\sqrt{2mE})
\right].
\]

A finite profile \(\eta_\epsilon(H-E)\) samples neighborhoods of both roots.
Its normalized expectation converges to their equal average. The equality of
the two sharp weights follows from the equal Jacobian magnitudes, not from an
extra probability postulate.

At \(E=0\), the two roots merge and \(H'(0)=0\). The regular-value theorem no
longer applies, and \(\Omega(E)\) diverges as \(E^{-1/2}\). This elementary
example shows why critical energies cannot be dismissed as a technical
footnote.

\section{Several simultaneous constraints}

For a smooth map
\[
C:M\to\mathbb R^r
\]
with full rank near \(C^{-1}(0)\), the normal Jacobian is
\[
\Jcal_C(z)
=
\sqrt{\det\!\bigl(DC(z)DC(z)^\ast\bigr)}.
\]
A normalized approximate identity
\(\eta_\epsilon^{(r)}\) in \(\mathbb R^r\) gives
\[
\int_M f(z)\eta_\epsilon^{(r)}(C(z))\,\dd\mu(z)
\longrightarrow
\int_{C^{-1}(0)}
\frac{f(z)}{\Jcal_C(z)}\,\dd\Sigma(z).
\]
This is the general coarea-tube result proved in the companion delta-kernel
paper. After dividing by the same expression with \(f=1\), one obtains a
normalized constraint-surface measure whenever its total mass is finite and
nonzero.

The full-rank hypothesis matters. First-class Hamiltonian constraints can
carry gauge redundancy, and singular constraint sets may require reduction
before a finite probability measure is meaningful. The finite-dimensional
coarea formula does not by itself construct an infinite-dimensional path
integral or a global gauge quotient.

\section{Canonical weighting is a different ensemble}

The canonical density
\[
\dd\nu_\beta
=
\frac{e^{-\beta H}}{Z(\beta)}\,\dd\mu
\]
weights many energies. It is not a finite-width approximation to one
microcanonical surface unless an additional concentration or thermodynamic
limit is proved.

Microcanonical and canonical ensembles can agree for suitable observables in
suitable large-system limits. They need not agree for finite systems, at
phase-transition points, or for systems with long-range interactions. The
shell theorem proves neither ensemble equivalence nor thermalization. It only
controls one regularized representation of a declared energy constraint.

\section{MTT interpretation and source obligation}

The mathematically complete downstream chain is
\[
\boxed{
\begin{gathered}
\text{regular Hamiltonian level geometry}\\
\downarrow\\
\text{normalized finite energy profile}\\
\downarrow\\
\text{coarea convolution and error bound}\\
\downarrow\\
\text{sharp microcanonical surface measure}.
\end{gathered}
}
\]

MTT can use this chain as an interface. An upstream construction would have to
emit:
\begin{enumerate}
\item the Hamiltonian or effective conserved quantity;
\item the prepared target energy \(E\);
\item the shell profile \(\eta\);
\item the physical width \(\epsilon\);
\item the phase-space measure and any reduction by constraints; and
\item the regime in which regularity and normalization hold.
\end{enumerate}

The common circle or fixed-point geometry may eventually help source a
conservation or return condition, but equality of interpretation is not a
derivation. No present theorem identifies \(\epsilon\) with a circle scale,
coherence length, spectral gap, or universal noise floor.

Different normalized profiles have the same sharp limit but different
finite-width predictions. That universality makes the limiting theorem robust,
while also showing why the limit cannot select the finite profile. Empirical
comparison at nonzero width would test the source model, not the coarea
formula.

\section{Status ledger}

\begin{center}
\begin{tabular}{p{0.20\linewidth}p{0.69\linewidth}}
\toprule
Status & Result\\
\midrule
Standard input &
Coarea representation of a regular energy-level distribution\\
Exact &
Normalized finite-shell measures converge to the sharp microcanonical measure\\
Exact &
The Lipschitz shell-bias bound separates numerator and normalization errors\\
Exact &
Every normalizable density depending only on \(H\) is Hamiltonian-flow
invariant\\
Exact example &
The one-dimensional kinetic shell gives equal weights at its two regular roots\\
Conditional &
Finite shell as detector resolution, preparation tolerance, or coarse graining\\
Open &
Critical-level theory, ergodicity, thermalization, and ensemble equivalence\\
Open in MTT &
Selected \(H\), \(E\), profile, width, and reduced phase-space measure\\
\bottomrule
\end{tabular}
\end{center}

\section{Conclusion}

The microcanonical delta is a precisely defined surface distribution. A
normalized finite energy shell is an ordinary probability measure that can
approach it. Coarea geometry turns the problem into convolution of
density-of-states functions, making both convergence and finite-width error
transparent.

The resulting picture is sharper than the slogan that the delta is merely a
hidden finite projection. At regular energies, many shell profiles share the
same limit; at finite width, they are distinct models. Hamiltonian flow
preserves all normalizable energy-profile shells exactly, but invariance alone
does not establish ergodicity or equilibrium.

For MTT, the theorem is now a reusable interface rather than a source claim.
Once selected geometry provides a conserved quantity, preparation energy, and
finite profile, the paper gives the normalized measure, sharp limit, and an
error certificate. Selecting those inputs remains the physical task.

\begin{thebibliography}{99}

\bibitem{EvansGariepy2015}
L.~C. Evans and R.~F. Gariepy,
\emph{Measure Theory and Fine Properties of Functions},
revised edition, CRC Press, Boca Raton, 2015,
doi:10.1201/b18333.

\bibitem{Ruelle1999}
D.~Ruelle,
\emph{Statistical Mechanics: Rigorous Results},
World Scientific, Singapore, 1999,
doi:10.1142/4090.

\bibitem{CampaDauxoisRuffo2009}
A.~Campa, T.~Dauxois, and S.~Ruffo,
\emph{Statistical Mechanics and Dynamics of Solvable Models with Long-Range
Interactions},
Physics Reports \textbf{480} (2009) 57--159,
doi:10.1016/j.physrep.2009.07.001.

\end{thebibliography}

\end{document}
